\chapterbackmatter{Supplementary Exercises}

\begin{enumerate}
\SupplementaryExercise{1}{The Transverse Delta Distribution and Coulomb Energy from Gauss's Law}
  {(Inspired by Daniel Harlow, Physics 8.324 Lecture Notes, Section 9.7 | Daniel Harlow, Physics 8.324 Lecture Notes, Sections 9.7 and 10.1.)}
  {harlow-qft2-electrodynamics-2024-ch8-curated-parent-1}
\begin{enumerate}
\item
Starting from the Coulomb-gauge mode expansion, derive the equal-time
commutator between \(A_i\) and its transverse momentum.  Show both its
momentum-space projector form and its position-space nonlocal form.
\item
Solve Gauss's law for \(A^0\) in Coulomb gauge with a localized charge
density.  Substitute the solution into the Hamiltonian and recover the
instantaneous bilocal Coulomb energy with its factor of one half.
\end{enumerate}

\SupplementaryExercise{2}{A Covariant Polarization Replacement}
  {(Inspired by Daniel Harlow, Physics 8.324 Lecture Notes, Section 11.5.)}
  {harlow-qft2-electrodynamics-2024-ch8-curated-parent-2}
For a photon attached to a conserved scattering current, prove that the
Coulomb-gauge polarization projector may be replaced by the covariant metric
inside a squared amplitude.  State exactly where current conservation enters.

\SupplementaryExercise{3}{Threshold Law for Pair Production and Tree-Level Electromagnetic \(R\) Ratio}
  {(Inspired by Daniel Harlow, Physics 8.324 Lecture Notes, Section 11.3 | Daniel Harlow, Physics 8.324 Lecture Notes, Section 11.4.)}
  {harlow-qft2-electrodynamics-2024-ch8-curated-parent-3}
\begin{enumerate}
\item
Without repeating the full annihilation calculation, use two-body phase
space and regularity of the tree amplitude to determine how the total
cross-section for producing a massive charged-fermion pair turns on at
threshold.
\item
At energies where fermion masses can be neglected, compare
\(e^+e^-\to\) hadrons with \(e^+e^-\to\mu^+\mu^-\).  Derive the
tree-level ratio in terms of accessible quark charges and color multiplicity.
\end{enumerate}

\SupplementaryExercise{4}{The Scalar-QED Ward Cancellation}
  {(Inspired by Daniel Harlow, Physics 8.324 Lecture Notes, Sections 10.1 and 11.5.)}
  {harlow-qft2-electrodynamics-2024-ch8-curated-parent-4}
Write the two exchange graphs and contact graph for scalar Compton
scattering.  Replace one photon polarization by its momentum and demonstrate
algebraically that all three contributions cancel.

\SupplementaryExercise{5}{Gauge Covariance, Current, and Vertices in Scalar QED}
  {(Inspired by David Tong, Lectures on Quantum Field Theory, Section 6.1 | David Tong, Lectures on Quantum Field Theory, Sections 6.1 and 6.5 | David Tong, Lectures on Quantum Field Theory, Section 6.5.1.)}
  {tong-qft-qed-2006-ch8-curated-parent-1}
\begin{enumerate}
\item
Determine the transformation of \(A_\mu\) that makes
\(D_\mu\Phi=(\partial_\mu-iqA_\mu)\Phi\) transform in the same way as
\(\Phi\).  Use it to prove gauge invariance of the scalar kinetic term.
\item
Vary the scalar-electrodynamics action with respect to \(A_\mu\) and derive
the full gauge-invariant matter current, including its term proportional to
\(A^\mu\).  Verify its conservation using the scalar field equation.
\item
Expand the scalar covariant-derivative term through second order in the
charge.  Derive the momentum-space rules for the photon--scalar vertex and
the two-photon seagull vertex, including their momentum dependence.
\end{enumerate}

\SupplementaryExercise{6}{Coulomb Exchange and Tree-Level Gauge Independence}
  {(Inspired by David Tong, Lectures on Quantum Field Theory, Section 6.6.2 | David Tong, Lectures on Quantum Field Theory, Section 6.6.3 | David Tong, Lectures on Quantum Field Theory, Sections 6.3--6.5.)}
  {tong-qft-qed-2006-ch8-curated-parent-2}
\begin{enumerate}
\item
Take the nonrelativistic limit of one-photon exchange between two charged
scalars.  Fourier transform the amplitude and show how equal and opposite
charges respectively produce repulsive and attractive Coulomb potentials.
\item
For the sum of the two tree diagrams in spinor Compton scattering, replace
the incoming photon polarization by its momentum.  Use the external Dirac
equations to show that the amplitude vanishes.
\item
Use a photon propagator with an arbitrary longitudinal gauge parameter.
Show that a tree amplitude connecting two conserved fermion currents is
independent of this parameter, and identify the assumption that would make
the conclusion fail.
\end{enumerate}

\SupplementaryExercise{7}{The Gupta--Bleuler Physical Space and Residual Gauge Modes}
  {(Adapted from Hong Liu and MIT OpenCourseWare, 8.323 Assignment 11, Problem 1.)}
  {mit-823-s23-ps11-ch8-curated-parent-1}
\begin{enumerate}
\item
In covariantly gauge-fixed free electrodynamics, calculate an equal-time
commutator of \(A_0\) with \(\partial_\mu A^\mu\).  Explain why its nonzero
value forbids imposing the Lorenz condition as an operator identity.
\item
Separate \(\partial_\mu A^\mu\) into annihilation and creation parts.
Formulate the Gupta--Bleuler condition, express it in terms of time-like and
longitudinal photon operators, and show why the vacuum satisfies it.
\item
Form the sum and difference of the time-like and longitudinal photon
annihilation operators.  Compute all four mixed commutators and identify
which creation operator generates the null excitations allowed by the
physical-state condition.
\item
Prove that every Gupta--Bleuler state is a transverse state plus a vector
orthogonal to every physical state.  Explain why quotienting by these null
vectors leaves a positive-definite two-polarization Hilbert space.
\item
Evaluate the one-particle wavefunction
\(\bra{\Omega}A_\mu(x)\ket{\mathbf k,\mathrm{null}}\).
Show that it is a pure gradient \(\partial_\mu\lambda\), with
\(\Box\lambda=0\), and interpret the result.
\item
Replace the null photon of definite momentum by a normalizable wave packet.
Construct the corresponding scalar gauge parameter explicitly and prove
that the complete wavefunction remains a residual Lorenz-gauge
transformation.
\end{enumerate}

\SupplementaryExercise{8}{Invariant Flux and Two-Body Phase Space}
  {(Inspired by Hong Liu and MIT OpenCourseWare, 8.323 Assignment 12, Notes 3--5 | Hong Liu and MIT OpenCourseWare, 8.323 Assignment 12, Notes 2--6 | Hong Liu and MIT OpenCourseWare, 8.323 Assignment 12, Notes 2--3.)}
  {mit-823-s23-ps12-ch8-curated-parent-1}
\begin{enumerate}
\item
Show that the flux denominator for two incident particles can be written as
\(4\sqrt{(p_1\cdot p_2)^2-m_1^2m_2^2}\).  Reduce it in the center-of-mass
frame and in the rest frame of the second particle.
\item
Starting from the Lorentz-invariant final-state measure, perform the energy
and radial momentum integrations for a general \(2\to2\) process.  Derive
the center-of-mass differential cross-section formula.
\item
Explain why integrating labeled momenta over the full final-state phase
space overcounts \(n\) identical particles by \(n!\).  Apply the argument
specifically to the two electrons in Møller scattering.
\end{enumerate}

\SupplementaryExercise{9}{Relativistic Rutherford Scattering}
  {(Adapted from Hong Liu and MIT OpenCourseWare, 8.323 Assignment 12, Problem 3.)}
  {mit-823-s23-ps12-ch8-curated-parent-2}
Treat a static Coulomb potential as an external field and compute
unpolarized elastic electron scattering in the Born approximation.  Derive
the relativistic spin correction multiplying the Rutherford result.

\SupplementaryExercise{10}{Electron--Muon Scattering Invariants and Heavy-Target Limit}
  {(Adapted from Hong Liu and MIT OpenCourseWare, 8.323 Assignment 12, Problem 4.)}
  {mit-823-s23-ps12-ch8-curated-parent-3}
\begin{enumerate}
\item
Calculate the tree amplitude for
\(e^-\mu^-\to e^-\mu^-\).  Sum and average over spins and express the
squared amplitude entirely in terms of \(s,t,u,m_e,m_\mu\).
\item
Convert the electron--muon result to the muon rest frame and then take
\(m_\mu\to\infty\) with the electron energy fixed.  Recover the
relativistic Rutherford formula and explain the disappearance of target
recoil.
\end{enumerate}

\SupplementaryExercise{11}{Electric Charge from a Static Potential and Dirac Magnetic Moment}
  {(Adapted from Bernhard Mistlberger and Kevin Zhou, Stanford PHYS 330 Problem Set 8, Problem 1.)}
  {zhou-stanford-phys330-ps8-2022-ch8-curated-parent-1}
\begin{enumerate}
\item
Evaluate the one-fermion matrix element of the QED interaction in a weak
static scalar potential.  Take the nonrelativistic limit and show that the
coefficient of \(A^0\) is the charge appearing in the covariant derivative.
\item
Use the Gordon identity to take the nonrelativistic limit of the spatial
Dirac current in a static magnetic field.  Read off the magnetic moment and
derive the tree-level gyromagnetic factor.
\end{enumerate}

\SupplementaryExercise{12}{A Pauli Term Shifts the Magnetic Moment and Anapole Coupling}
  {(Adapted from Bernhard Mistlberger and Kevin Zhou, Stanford PHYS 330 Problem Set 8, Problem 4.)}
  {zhou-stanford-phys330-ps8-2022-ch8-curated-parent-2}
\begin{enumerate}
\item
Add the gauge-invariant dimension-five Pauli interaction
\(\bar\psi\sigma^{\mu\nu}\psi F_{\mu\nu}\).  Reduce it for slowly moving
fermions and determine how its coefficient shifts the magnetic moment.
\item
Reduce
\(\bar\psi\sigma^{\mu\nu}\gamma_5\psi F_{\mu\nu}\) in the
nonrelativistic limit.  Identify the induced dipole interaction and
determine its transformation under parity and time reversal.
\item
Analyze the dimension-six operator
\(\bar\psi\gamma^\mu\psi\,\partial^\nu F_{\mu\nu}\).
Use Maxwell's equation to relate it to a contact interaction with an
external conserved current, and describe its form-factor interpretation.
\item
Study
\(\bar\psi\gamma^\mu\gamma_5\psi\,\partial^\nu F_{\mu\nu}\).
Find its nonrelativistic coupling to an external current, explain why it
does not describe an ordinary electric or magnetic dipole, and classify
its discrete symmetries.
\end{enumerate}
\end{enumerate}
